\chapter{命令通路与运行时控制}
\oneline{上位机经 UDP 6001 广播命令,cvp\_cmd\_rx 在 RGMII 域极简解析后送入
100M 域命令寄存器;手动曝光/模拟增益经 I2C 请求/确认桥写入传感器,AE 则调节
传感器曝光与 FPGA 数字增益;白平衡增益即时覆盖 color\_correct,状态经 UDP 6002
回传。本章讲命令基础设施与手动/自动
切换语义——AE/AWB 闭环的算法本体见第 3 章。}

\section{命令通路:cvp\_cmd\_rx}

cvp\_cmd\_rx 直接在 RGMII 接收域做极简 UDP 命令解析,不经过协议栈:IDDR 还原
字节流 $\to$ 过滤状态机逐条校验(SFD、EtherType=IPv4、协议=UDP、目的端口 6001、
魔数 C5A55001,任一不符即丢帧)$\to$ 命令事件跨时钟域 $\to$ 100\,MHz 域命令寄存器
组。设计上有意不校验 MAC/IP 地址(命令走广播),宽容上位机与脚本混用。

命令域覆盖:曝光码(半行编码)、模拟增益档、白平衡三通道增益(Q8)与模式
(预设/色温 K/手动/自动)、R/B 互换、算法模式与阈值参数。手动曝光/模拟增益经
I2C 桥重写传感器(约 1$\sim$2 帧生效);白平衡增益直接覆盖 color\_correct
(即拖即生效)。自动曝光由 PID 控制器优先调整传感器曝光,到达上限后调节 ISP
中的 Q8 数字增益;模拟增益保持用户设定。状态与应答包经 UDP 6002
回传(含 CRC 与 IPv4 校验和)。

\section{I2C 请求/确认桥:传感器曝光执行通道}

传感器寄存器写入(上电初始化、AE 曝光更新与手动设置)共用单一的请求/确认桥:任何时刻
只有一个在飞的写事务,写入中的数据保持不变,I2C 异常以 NACK/超时上报到状态口。
这一设计同时服务初始化和运行期写入——初始化按配置表(主表/AE 表/增益表)批量写入,
运行期由手动命令更新传感器曝光码与模拟增益档,AE 闭环(第 3 章)则通过该桥更新
曝光码,并在像素处理域独立锁存数字增益;自动控制不改写模拟增益档。摄像头启动流程带
等待与重试,并以第一完整帧门控后续流水,上电即出流无需手动干预。

\section{手动/自动切换语义}

GUI 与命令行的操作语义固定:拖动曝光或模拟增益即关闭自动曝光转入手动;拖动 RGB
增益即切换手动白平衡;自动模式下的参数写入同样走确认应答。GUI 显示的是确认写入
的值——界面不声称读回了传感器寄存器,传感器生效可能延后数帧。边缘/素描的阈值
参数与算法模式切换同属命令域,切换在帧首生效、无撕裂。

\section{状态与诊断}

运行状态经两个通道可观测:UDP 6002 状态/应答包(GUI 实时显示写入确认、I2C 健康、
帧计数与流错误;v29 起并承担链路健康判定——GUI 能收图即全链路正常,收不到图先排查
PC 侧 6000 端口占用)与 GUI 状态栏(fps 与完整率,正常 100\%)。测量纪律:UDP 抓帧
必须独占 6000 端口(只打开一个收图程序),否则端口分流会产生假读数。
